\section*{Hoofdstuk \ref{ch:b1:intermolecular-forces-solvents} --- Intermoleculaire krachten en oplosmiddelen}

\begin{solution}{exo:b1:intermolecular-forces-solvents:1}
Argon: alleen London. Waterstofchloride: Keesom, Debye en London, waarbij
London overheerst (\cref{exo:b1:intermolecular-forces-solvents:8});
chloor is te groot voor sterke \omterm{def:b1:intermolecular-forces-solvents:hbond}{waterstofbruggen}. Water: \omterm{def:b1:intermolecular-forces-solvents:hbond}{waterstofbruggen}
overheersen, naast de drie vanderwaalstermen. Methaan: alleen London.
Dijood: alleen London, maar sterk omdat het molecuul groot en zeer
polariseerbaar is.
\end{solution}

\begin{solution}{exo:b1:intermolecular-forces-solvents:2}
Tussen hun eigen moleculen: \ce{CH3OH} (\ce{O-H}), \ce{CH3NH2}
(\ce{N-H}) en \ce{HF}. \ce{CH3OCH3} en \ce{CH3F} hebben \omterm{def:b1:lewis-resonance-vsepr:lewis}{vrije paren} op O
of F maar geen waterstof daarop: ze kunnen alleen \omterm{def:b1:intermolecular-forces-solvents:hbond}{waterstofbruggen}
opnemen, bijvoorbeeld van water. \ce{CH4} geeft noch neemt op.
\end{solution}

\begin{solution}{exo:b1:intermolecular-forces-solvents:3}
Water en ethanol: polair, protisch. Propanon en acetonitril: polair,
aprotisch. Dichloormethaan: zwak polair, aprotisch. Hexaan: apolair,
aprotisch.
\end{solution}

\begin{solution}{exo:b1:intermolecular-forces-solvents:4}
Dijood is apolair: in hexaan ruilt het \omterm{def:b1:intermolecular-forces-solvents:vdw}{londoninteracties} in voor
\omterm{def:b1:intermolecular-forces-solvents:vdw}{londoninteracties}, terwijl het in water \omterm{def:b1:intermolecular-forces-solvents:hbond}{waterstofbruggen} zou verbreken
zonder ze te vervangen. Natriumchloride is ionair: alleen een
oplosmiddel met een hoge permittiviteit kan zijn ionen scheiden en
alleen een \omterm{def:b1:intermolecular-forces-solvents:solvent-classes}{polair oplosmiddel} kan ze solvateren; hexaan
($\varepsilon_r = 1.89$) doet geen van beide.
\end{solution}

\begin{solution}{exo:b1:intermolecular-forces-solvents:5}
Methoxymethaan heeft geen \ce{O-H}: het kan geen \omterm{def:b1:intermolecular-forces-solvents:hbond}{waterstofbruggen} geven
en kookt het laagst. Methanol en ethanol vormen allebei \omterm{def:b1:intermolecular-forces-solvents:hbond}{waterstofbruggen};
ethanol, met één \ce{CH2} meer, heeft sterkere \omterm{def:b1:intermolecular-forces-solvents:vdw}{londoninteracties}.
Ethaanzuur vormt dimeren via \omterm{def:b1:intermolecular-forces-solvents:hbond}{waterstofbruggen} en is het zwaarst: het
kookt het hoogst.
\end{solution}

\begin{solution}{exo:b1:intermolecular-forces-solvents:6}
In vacuüm is $F = e^2/(4\pi\varepsilon_0 r^2) = (1.602 \times
10^{-19})^2/(4\pi \times 8.854 \times 10^{-12} \times (5.00 \times
10^{-10})^2) = \qty{9.23e-10}{N}$. In water is $F/78.4 = \qty{1.18e-11}{N}$;
in hexaan $F/1.89 = \qty{4.88e-10}{N}$. Water verzwakt de aantrekking
41 keer meer dan hexaan: water kan de ionen uit elkaar houden, hexaan
niet.
\end{solution}

\begin{solution}{exo:b1:intermolecular-forces-solvents:7}
Mengen van propanon met water verbreekt een deel van de \omterm{def:b1:intermolecular-forces-solvents:hbond}{waterstofbruggen}
tussen watermoleculen, maar vormt er nieuwe tussen water en de zuurstof
van \ce{C=O}, plus dipool-dipoolinteracties: de balans is gunstig.
Hexaan biedt water alleen zwakke \omterm{def:b1:intermolecular-forces-solvents:vdw}{londoninteracties} in ruil voor de
\omterm{def:b1:intermolecular-forces-solvents:hbond}{waterstofbruggen} die het zou verbreken: de twee vloeistoffen blijven
gescheiden.
\end{solution}

\begin{solution}{exo:b1:intermolecular-forces-solvents:8}
\ce{HBr} heeft meer elektronen en een grotere, polariseerbaardere wolk:
zijn \omterm{def:b1:intermolecular-forces-solvents:vdw}{londoninteractie} overtreft die van \ce{HCl} met meer dan zijn
\omterm{def:b1:intermolecular-forces-solvents:vdw}{keesominteractie} tekortschiet. London overheerst.
\end{solution}

\begin{solution}{exo:b1:intermolecular-forces-solvents:9}
\chemfig{CH_3-C(=[:60]O)-[:-60]O-[:0]H} tegenover zijn partner, met twee
bindingen $\ce{O-H}\cdots\ce{O=C}$ die een achtring sluiten. In benzeen, apolair en aprotisch, concurreert niets
met deze bindingen: het zuur is gedimeriseerd en gedraagt zich als een
deeltje van ongeveer \qty{120}{g/mol}. In water vormt elk zuurmolecuul in
plaats daarvan \omterm{def:b1:intermolecular-forces-solvents:hbond}{waterstofbruggen} met water: de dimeren vallen uiteen.
\end{solution}

\begin{solution}{exo:b1:intermolecular-forces-solvents:10}
Natriumchloride bestaat uit ionen: water, met een hoge permittiviteit,
scheidt ze (dissociatie) en hydrateert ze (zuurstof naar \ce{Na+},
waterstof naar \ce{Cl-}). Waterstofchloride is een \omterm{def:b1:lewis-resonance-vsepr:dipole}{polair molecuul}:
water, dat polair is, zet de \ce{H-Cl}-binding om in een ionenpaar
\ce{H3O+ Cl-} (ionisatie), scheidt het paar (dissociatie) en hydrateert de
ionen. Hexaan is apolair en heeft een lage permittiviteit: het ioniseert
noch dissocieert, zodat er geen ionen zijn om de stroom te geleiden.
\end{solution}

\begin{solution}{exo:b1:intermolecular-forces-solvents:11}
$(174.1 - 36.1)/5 = \qty{27.6}{\celsius}$ per \ce{CH2}. De laatste
toename (nonaan naar decaan) is \qty{23.4}{\celsius}, dus undecaan zou
rond $174 + 22 \approx \qty{196}{\celsius}$ moeten koken. Elke \ce{CH2}
voegt dezelfde londonbijdrage toe, maar een steeds kleiner deel van het
totaal van het molecuul, terwijl de temperatuur die nodig is om de
interacties te overwinnen al hoog is.
% ledger: bp:C9H20
\end{solution}

\begin{solution}{exo:b1:intermolecular-forces-solvents:12}
\omterm{def:b1:intermolecular-forces-solvents:amphiphile}{Hydrofiel} deel: de sulfaatkop \ce{-OSO3^-} met zijn tegenion \ce{Na+};
\omterm{def:b1:intermolecular-forces-solvents:amphiphile}{hydrofoob} deel: de keten van twaalf koolstofatomen. Aan het oppervlak de
koppen in het water en de staarten in de lucht; in een micel de staarten
binnen en de koppen buiten. De staarten lossen op in het vet van de
vlek, de koppen houden het in contact met water, en het vet wordt in
micellen losgemaakt.
\end{solution}

\begin{solution}{pb:b1:intermolecular-forces-solvents:1}
\textbf{1.} De \omterm{def:b1:intermolecular-forces-solvents:vdw}{londoninteractie}, de enige tussen apolaire atomen.
\textbf{2.} Van helium tot xenon groeien het aantal elektronen en de
grootte, en dus de \omterm{def:b1:periodicity:polarisability}{polariseerbaarheid} en de londonaantrekking.
\textbf{3.} Helium heeft twee elektronen die stevig door zijn kern
worden vastgehouden: het minst polariseerbare atoom, met de zwakste
londonaantrekking.
\textbf{4.} $68.8 - 36.1 = \qty{32.7}{\celsius}$: één \ce{CH2} voegt
elektronen en contactoppervlak toe, en dus londonaantrekking.
\textbf{5.} London, dat groeit met de \omterm{def:b1:periodicity:polarisability}{polariseerbaarheid}:
$\ce{CH4} < \ce{SiH4} <
\ce{GeH4}$.
\textbf{6.} Koolstof is niet elektronegatief genoeg om \ce{C-H} sterk te
polariseren, en methaan heeft geen \omterm{def:b1:lewis-resonance-vsepr:lewis}{vrij paar} om een \omterm{def:b1:intermolecular-forces-solvents:hbond}{waterstofbrug} op te
nemen.
\textbf{7.} Water: polair, protisch, met de hoogste permittiviteit; het
scheidt en solvateert beide ionen.
\textbf{8.} $78.4/1.89 = 41$: de aantrekking is in water 41 keer zwakker.
\textbf{9.} Acetonitril (of propanon): polair genoeg om het zout op te
lossen, aprotisch zodat het anion vrij blijft.
\textbf{10.} Hexaan: niet mengbaar met water, apolair zoals dijood. Het
jood gaat over naar het hexaan, de bovenste laag (hexaan is minder dicht
dan water).
\textbf{11.} Ethanol is \omterm{def:b1:intermolecular-forces-solvents:miscibility}{mengbaar} met water: er vormt zich geen tweede
laag.
\textbf{12.} Zijn zuurstofatoom neemt \omterm{def:b1:intermolecular-forces-solvents:hbond}{waterstofbruggen} van water op, maar
het kan er zelf geen geven en zijn twee ethylgroepen zijn \omterm{def:b1:intermolecular-forces-solvents:amphiphile}{hydrofoob}: het
lost maar weinig op.
\textbf{13.} $-66.4 - (-85.0) = \qty{18.6}{\celsius}$ per \omterm{def:b1:periodicity:table}{periode};
geëxtrapoleerd \ce{HF}: $-85.0 - 18.6 = \qty{-103.6}{\celsius}$.
\textbf{14.} \ce{HF} kookt bij \qty{19.6}{\celsius}, \qty{123}{\celsius}
boven de extrapolatie.
\textbf{15.} $-62.5 - (-87.8) = \qty{25.3}{\celsius}$; geëxtrapoleerd
\ce{NH3}: \qty{-113.1}{\celsius}; werkelijk $-33.4$, \qty{80}{\celsius}
hoger.
\textbf{16.} $-88.1 - (-112.0) = \qty{23.9}{\celsius}$; geëxtrapoleerd
\ce{CH4}: \qty{-135.9}{\celsius}; werkelijk $-161.5$, eronder. Methaan
vormt geen \omterm{def:b1:intermolecular-forces-solvents:hbond}{waterstofbruggen}; het is gewoon klein en weinig
polariseerbaar.
\textbf{17.} Water (ongeveer \qty{179}{\celsius}) $>$ \ce{HF} (123) $>$
\ce{NH3} (80).
\textbf{18.} \ce{NH3} kan er drie geven maar slechts één opnemen (één
\omterm{def:b1:lewis-resonance-vsepr:lewis}{vrij paar}); \ce{HF} kan er drie opnemen maar slechts één geven. Water
geeft er twee en neemt er twee op: elk molecuul kan deelnemen aan vier
\omterm{def:b1:intermolecular-forces-solvents:hbond}{waterstofbruggen}, een volledig driedimensionaal netwerk.
\textbf{19.} \ce{H2S} en \ce{H2Se} vormen geen noemenswaardige
\omterm{def:b1:intermolecular-forces-solvents:hbond}{waterstofbruggen}: hun kookpunten volgen de londontrend die water zonder
\omterm{def:b1:intermolecular-forces-solvents:hbond}{waterstofbruggen} zou volgen. Water zelf is juist de afwijking die gemeten
wordt.
\textbf{20.} $-41.3 - (-60.3) = \qty{19.0}{\celsius}$ per \omterm{def:b1:periodicity:table}{periode}.
\textbf{21.} $T(p) = -60.3 + 19.0\,(p - 3)$ in \unit{\celsius}.
\textbf{22.} $T(2) = -60.3 - 19.0 = \qty{-79.3}{\celsius}$.
\textbf{23.} $100.0 - (-79.3) = \qty{179}{\celsius}$.
\textbf{24.} Zonder zijn \omterm{def:b1:intermolecular-forces-solvents:hbond}{waterstofbruggen} zou water koken bij ongeveer
\textbf{\qty{-79}{\celsius}}, en op aarde een gas zijn.
\end{solution}
